% ==========================================================================
% Capítulo: Objetivos
% Estado: se registra en ../STATUS.md, no aquí.
%
% Debe contener (project-context/formato-anteproyecto.md, "Objetivos";
% extensión máxima 1 página):
%   - OBJETIVO GENERAL: uno solo, verbo en infinitivo, medible, alcanzable
%     en ~8 meses, en relación directa con el título del proyecto.
%   - OBJETIVOS ESPECÍFICOS: máximo 4, derivados del general y que en
%     conjunto lo cubran; precisos, realizables y mensurables.
% NO debe contener: actividades disfrazadas de objetivos, verbos
% compuestos, ni enunciados no medibles.
%
% Reglas de formulación y anti-patrones: ../../skills/objectives-writting/
% (objectives-definition.md y objectives-avoid.md) — de obligada consulta.
% Objetivos base a reformular: ../../project-context/documentation.md,
% sección "Formulation of Objectives".
% ==========================================================================
\chapter{Objetivos}\label{cha:objetivos}

\section{Objetivo general}

Construir la plataforma de orquestación y gobernanza de cargas de inteligencia artificial y aprendizaje automático del IAsLab, mediante módulos de despliegue de inferencia, gobernanza de recursos por rol y observabilidad de infraestructura, para consolidar el ciclo de vida de MLOps \emph{on-premise} en la Universidad Icesi.

\section{Objetivos específicos}

\begin{enumerate}
  % Hace que \cref/\Cref nombren estos ítems como "objetivo específico N"
  % en lugar de "Punto N"; es local a este entorno (ver ../../skills/latex/preamble.tex).
  \crefalias{enumi}{objesp}
  \item \label{obj:telemetria} Desarrollar un módulo de observabilidad de la plataforma para la captura y visualización de métricas de \emph{hardware} y de carga de trabajo de los nodos de cómputo del IAsLab, orientado a la supervisión del estado de dichos nodos.

  \item \label{obj:gobernanza} Implementar un sistema de cuotas y restricciones de infraestructura, acoplado a los roles provistos por SAAMFI y a un rol de administrador de la plataforma, que reserve cupos de cómputo para las cargas de inferencia con prioridad de los profesores sobre la asignación de la capacidad de la sala, con el fin de prevenir la monopolización de las GPU.

  \item \label{obj:orquestacion} Implementar sobre el clúster de cómputo del IAsLab el despliegue concurrente de modelos de inteligencia artificial ya entrenados, mediante motores de inferencia empaquetados como imágenes de contenedor, con el fin de llevar los modelos del laboratorio a una fase de servicio consumible.

  \item \label{obj:evaluacion} Evaluar el desempeño y la aceptación de la plataforma mediante pruebas de carga con veinte usuarios simultáneos y la aplicación del \emph{System Usability Scale} a usuarios del laboratorio, con el fin de determinar si la plataforma sostiene esa concurrencia y en qué banda de referencia publicada se ubica su usabilidad percibida.
\end{enumerate}
